% tokens.tex — improvement #79
%
% Every value the print design can be tuned by, in one file. Nothing else under
% scripts/tex/ hardcodes a length, a size or a grey: they all read the names defined here.
%
% This is the file improvement #77 edits during calibration. If a change to the look of the
% book needs an edit anywhere else, the token for it is missing.
%
% ---------------------------------------------------------------------------
% Palette — five tones, chosen for how ink behaves rather than how a screen looks
% ---------------------------------------------------------------------------
%
% The book prints in black and white, so the whole design is carried by rule weight, type
% weight and structure, with two tints and no more. Tints under 6% vanish on uncoated
% stock, and anything from 13% up goes muddy behind 10pt text — which is why there is no
% third step, and why #81 distinguishes the three callout types by structure instead of
% by fill.

\usepackage{xcolor}

\definecolor{ink}{gray}{0.00}       % body, chapter titles, H2, H3
\definecolor{inkmid}{gray}{0.45}    % eyebrow, captions, folio, running head — 55% K
\definecolor{inklight}{gray}{0.65}  % hairlines, table rules — 35% K
\definecolor{tintone}{gray}{0.94}   % code ground, table zebra — 6% K
\definecolor{tinttwo}{gray}{0.88}   % warning-callout ground — 12% K

% `gray` in xcolor is *lightness*, so a 55% K ink is gray 0.45. Getting this backwards
% produces a book that looks correct on screen and prints inverted in weight.

% ---------------------------------------------------------------------------
% Type scale
% ---------------------------------------------------------------------------
%
% Starting values. #77 measures the real lines-per-page rate and tunes these against the
% 700-page ceiling in BOOK-SPEC decision #13.

% 🔴 Calibrated at #87, and the calibration is mostly *not* here. The slack a page has is
% vertical and horizontal, and on A4 only one of them can be spent. The text block is
% 156mm at 10pt, which is already about 85 characters a line — above the 66–80 a reader
% tracks comfortably. Widening the measure or shrinking the type would buy pages by
% making the book harder to read, so neither moved. What moved is leading, paragraph
% spacing, block spacing and the top and bottom margins. See BOOK-SPEC decision #19.
\newcommand{\tokBodySize}{10pt}
\newcommand{\tokBodyLeading}{12.8pt}   % 1.28 — #87. Tight for an 85-character measure, and the floor for one
\newcommand{\tokLeadRatio}{1.067}      % \tokBodyLeading ÷ the 12pt the 10pt class would use
\newcommand{\tokSmallSize}{8.5pt}      % eyebrows, labels, table body
\newcommand{\tokMicroSize}{8pt}        % running heads and folios
\newcommand{\tokCodeSize}{8.5pt}
\newcommand{\tokCodeLeading}{10pt}      % #87 — 1.18 at 8.5pt, still open for a monospace
\newcommand{\tokDeckSize}{11pt}        % the one-sentence promise under a chapter title
\newcommand{\tokChapterSize}{22pt}
\newcommand{\tokPartNumberSize}{60pt}
\newcommand{\tokPartTitleSize}{28pt}

% ---------------------------------------------------------------------------
% Page geometry — consumed by structure.tex at #80
% ---------------------------------------------------------------------------
%
% Mirrored, because a uniform margin puts the gutter and the fore-edge at the same width
% and the text block then drifts toward the spine on every recto.

\newcommand{\tokPaper}{a4paper}
\newcommand{\tokMarginInner}{26mm}
\newcommand{\tokMarginOuter}{28mm}
% #87 took 2mm off the head and 3mm off the foot — 5mm of text height, about one extra
% line a page. The inner and outer margins did **not** move: they set the measure, and the
% measure is already at the limit of what reads well.
\newcommand{\tokMarginTop}{20mm}
\newcommand{\tokMarginBottom}{21mm}

% ---------------------------------------------------------------------------
% Rules and spacing
% ---------------------------------------------------------------------------

\newcommand{\tokHairline}{0.4pt}       % table rules, callout boxes
\newcommand{\tokAccentRule}{1.2pt}     % under a chapter title
\newcommand{\tokBarRule}{2.5pt}        % the left bar on a gotcha callout
\newcommand{\tokTracking}{140}         % +140/1000 em on small caps labels
\newcommand{\tokParSkip}{0.22em}       % 0.9em is a web value; on paper it costs ~90 pages. 0.30 → 0.22 at #87
\newcommand{\tokCellPad}{4pt}
\newcommand{\tokCodePad}{5pt}
\newcommand{\tokCalloutPad}{7pt}

% What it costs to break a line inside an inline code span — #77. Pandoc sets a code span
% as one unbreakable run, which is correct in a 442pt measure and ruinous in a 53pt table
% column: it was the single cause of every overfull box in the book wider than 20pt.
% scripts/lua/callouts.lua offers TeX a break after each path, dot, colon and camelCase
% boundary inside a span; this is how badly it should want to avoid taking one. 500 is
% high enough that a code span in body prose still sets in one piece, and low enough that
% a narrow cell breaks rather than bleeding into the gutter.
\newcommand{\tokCodeBreakPenalty}{500}

% And what a break *anywhere else* in a span costs — offered only inside a table cell,
% where the column may be narrower than the shortest piece the rule above can produce.
% Four times the cost, so it stays the last resort rather than the first fit.
\newcommand{\tokCodeBreakPenaltyAny}{2000}

% ---------------------------------------------------------------------------
% Blocks — consumed by blocks.tex at #81
% ---------------------------------------------------------------------------
%
% The book prints in black and white, so the three callout types are told apart by
% structure rather than by fill: a full box, a pair of rules, and a solid left bar over
% \tinttwo. Greyscale collapses the reference design's three hues to within one percent
% of each other; rule weight survives greyscale, photocopying and e-ink.

\newcommand{\tokTableSize}{9pt}        % table body — one step under the body serif
\newcommand{\tokTableLeading}{11pt}    % #87
\newcommand{\tokPullSize}{12pt}        % a line lifted out of the prose
\newcommand{\tokPullLeading}{15pt}
\newcommand{\tokPullClear}{10pt}       % air above and below a pull quote
\newcommand{\tokPillPad}{4pt}
\newcommand{\tokBlockSkip}{7.5pt}      % air around every block in this library — 9pt → 7.5pt at #87

% 🔴 The same value as a *length register*, and blocks.tex must use this one wherever it
% scales the skip. `-0.4\tokBlockSkip` expands to `-0.4 7.5pt`, which is not a dimension
% and halts the build. Before #87 the token was `9pt`, so the same expression silently
% read as **-0.49pt** — a plausible-looking number that was never the 3.6pt the code
% meant. A macro that holds a dimension cannot be multiplied; a register can.
\newlength{\bookBlockSkip}
\setlength{\bookBlockSkip}{\tokBlockSkip}

% ---------------------------------------------------------------------------
% Diagrams — consumed by blocks.tex at #82
% ---------------------------------------------------------------------------
%
% The 96 Mermaid figures are placed at their natural size and only ever shrunk, so this
% is a ceiling rather than a size. A diagram given the whole text height fills the page
% alone and pushes the sentence that reads it over the fold; 0.72 leaves prose on both
% sides of every figure in the book.

\newcommand{\tokDiagramMaxHeight}{0.72\textheight}

% The size the *labels inside* a diagram should end up at once it has been fitted. Read
% by scripts/lua/mermaid.lua, not by any TeX in this repository, which is why it carries
% its unit as a bare number: Lua has to do arithmetic with it.
%
% It is a target rather than a setting. A Mermaid diagram is rendered at its own natural
% size and then shrunk to fit the measure, so the type size on the page is a function of
% how wide the diagram happens to be — the 96 figures in this book shrink by anything
% from 1.00 to 0.36. The filter renders each one twice: once to find out how far it will
% shrink, and again at a font size chosen so that the shrink lands here. 8.5pt is the
% same step the design already uses for table body, eyebrows and callout labels.
\newcommand{\tokDiagramLabelSize}{8.5}
